% from ega2-1-fr.tex
\oldpage[II]{79}

  il suffit qu'il existe un entier $n\geq0$, un sous-$\mathcal{O}_Y$-Module
  quasi-cohérent de type fini $\mathcal{E}$ de $\mathcal{S}_n$, tels que
  l'homomorphisme
  \[
    \psi'=\psi_n\circ q^*(j):q^*(\mathcal{E})
      \longrightarrow\mathcal{L}^{\otimes n}=\mathcal{L}'
  \]
  ($j$ étant l'injection $\mathcal{E}\to\mathcal{S}_n$) soit surjectif et que
  le morphisme
  $r_{\mathcal{L}',\psi'}:X\to\mathbf{P}(\mathcal{E})$ soit une immersion.
  \item[\rm{(ii)}] Soient $\mathcal{F}$ un $\mathcal{O}_Y$-Module
  quasi-cohérent, $\varphi:q^*(\mathcal{F})\to\mathcal{L}$ un homomorphisme
  surjectif. Pour que le morphisme $r_{\mathcal{L},\varphi}$ soit une
  immersion $X\to\mathbf{P}(\mathcal{F})$, il faut et il suffit qu'il existe
  un sous-$\mathcal{O}_Y$-Module quasi-cohérent de type fini $\mathcal{E}$ de
  $\mathcal{F}$ tel que l'homomorphisme
  $\varphi'=\varphi\circ q^*(j):q^*(\mathcal{E})\to\mathcal{L}$ (où
  $j:\mathcal{E}\to\mathcal{F}$ est l'injection canonique) soit surjectif et
  que $r_{\mathcal{L},\varphi'}:X\to\mathbf{P}(\mathcal{E})$ soit une
  immersion.
\end{enumerate}
\end{proposition}

\begin{proof}
\begin{enumerate}
  \item[\rm{(i)}] Le fait que $r_{\mathcal{L},\psi}$ soit partout défini et
  soit une immersion équivaut d'après
  (\hyperref[II.3.8.5-fr]{3.8.5}) à l'existence de $n\geq0$ et de
  $\mathcal{E}$ tels que, si $\mathcal{S}'$ est la sous-algèbre de
  $\mathcal{S}$ engendrée par $\mathcal{E}$, l'homomorphisme
  $q^*(\mathcal{E})\to\mathcal{L}^{\otimes n}$ soit surjectif et que
  $r_{\mathcal{L}',\psi'}:X\to\operatorname{Proj}(\mathcal{S}')$ soit partout
  défini et soit une immersion. On a d'autre part un homomorphisme canonique
  surjectif $\mathbf{S}(\mathcal{E})\to\mathcal{S}'$ auquel correspond une
  immersion fermée
  $\operatorname{Proj}(\mathcal{S}')\to\mathbf{P}(\mathcal{E})$
  (\hyperref[II.3.6.2-fr]{3.6.2})~; d'où la conclusion.
  \item[\rm{(ii)}] Comme $\mathcal{F}$ est limite inductive de ses
  sous-Modules quasi-cohérents de type fini $\mathcal{E}_\lambda$
  (\hyperref[I.9.4.9-fr]{I, 9.4.9}), $\mathbf{S}(\mathcal{F})$ est limite
  inductive des $\mathbf{S}(\mathcal{E}_\lambda)$~; la conclusion résulte
  alors de (\hyperref[II.3.8.4-fr]{3.8.4}), en observant que, dans la
  démonstration de (\hyperref[II.3.8.4-fr]{3.8.4}), on peut prendre tous les
  $n_i$ égaux à $1$~: en effet ($Y$ étant supposé affine), si
  $r=r_{\mathcal{L},\varphi}$ est une immersion, $r(X)$ est un sous-espace
  quasi-compact de $\mathbf{P}(\mathcal{F})$ que l'on peut recouvrir par un
  nombre fini d'ouverts de $\mathbf{P}(\mathcal{F})$ de la forme $D_+(f)$ avec
  $f\in F$, tels que $D_+(f)\cap r(X)$ soit fermé.
\end{enumerate}
\end{proof}

\begin{definition}[4.4.2]
\label{II.4.4.2-fr}
Soient $Y$ un préschéma, $q:X\to Y$ un morphisme. On dit qu'un
$\mathcal{O}_X$-Module inversible $\mathcal{L}$ est \emph{très ample pour
$q$} (ou \emph{très ample pour $Y$}, ou simplement \emph{très ample} si
aucune confusion n'en résulte) s'il existe un $\mathcal{O}_Y$-Module
quasi-cohérent $\mathcal{E}$ et une $Y$-immersion $i$ de $X$ dans
$P=\mathbf{P}(\mathcal{E})$ tels que $\mathcal{L}$ soit isomorphe à
$i^*(\mathcal{O}_P(1))$.
\end{definition}

Il revient au même (\hyperref[II.4.2.3-fr]{4.2.3}) de dire qu'il existe un
$\mathcal{O}_Y$-Module quasi-cohérent $\mathcal{E}$ et un homomorphisme
surjectif $\varphi:q^*(\mathcal{E})\to\mathcal{L}$ tels que
$r_{\mathcal{L},\varphi}:X\to\mathbf{P}(\mathcal{E})$ soit une immersion.

On notera que l'existence d'un $\mathcal{O}_X$-Module très ample relativement
à $Y$ implique que $q$ est séparé
((\hyperref[II.3.1.3-fr]{3.1.3}) et
(\hyperref[I.5.5.1-fr]{I, 5.5.1, (i) et (ii)})).

\begin{corollary}[4.4.3]
\label{II.4.4.3-fr}
Supposons qu'il existe une $\mathcal{O}_Y$-Algèbre graduée quasi-cohérente
$\mathcal{S}$ engendrée par $\mathcal{S}_1$, et une $Y$-immersion
$i:X\to P=\operatorname{Proj}(\mathcal{S})$ telle que $\mathcal{L}$ soit
isomorphe à $i^*(\mathcal{O}_P(1))$~; alors $\mathcal{L}$ est très ample
relativement à $q$.
\end{corollary}

\begin{proof}
En effet, si $\mathcal{F}=\mathcal{S}_1$, l'homomorphisme canonique
$\mathbf{S}(\mathcal{F})\to\mathcal{S}$ est surjectif, donc, en composant
l'immersion fermée correspondante
$\operatorname{Proj}(\mathcal{S})\to\mathbf{P}(\mathcal{F})$
(\hyperref[II.3.6.2-fr]{3.6.2}) et l'immersion $i$, on obtient une immersion
$j:X\to\mathbf{P}(\mathcal{F})=P'$ telle que $\mathcal{L}$ soit isomorphe à
$j^*(\mathcal{O}_{P'}(1))$ (\hyperref[II.3.6.3-fr]{3.6.3}).
\end{proof}

\begin{proposition}[4.4.4]
\label{II.4.4.4-fr}
Soient $q:X\to Y$ un morphisme quasi-compact, $\mathcal{L}$ un
$\mathcal{O}_X$-Module inversible. Les propriétés suivantes sont
équivalentes~:
\begin{enumerate}
  \item[\rm{a)}] $\mathcal{L}$ est très ample relativement à $q$.
  \item[\rm{b)}] $q_*(\mathcal{L})$ est quasi-cohérent, l'homomorphisme
  canonique $\sigma:q^*(q_*(\mathcal{L}))\to\mathcal{L}$ est surjectif, et le
  morphisme
  $r_{\mathcal{L},\sigma}:X\to\mathbf{P}(q_*(\mathcal{L}))$ est une immersion.
\end{enumerate}
\end{proposition}

\begin{proof}
Comme $q$ est quasi-compact, on sait que $q_*(\mathcal{L})$ est quasi-cohérent
si $q$ est séparé (\hyperref[I.9.2.2-fr]{I, 9.2.2}).

